% -----------------------------------------------------------------------------
%  entries/q023.tex
%  Issued 14/10/2026. ID suggested by Claude; Gabriel's own numbering is the
%  authority. Built from the 50%-rule step of the research brief
%  (claude/research/cv-hypergraph-negativity/brief.md, section 2(b)).
%  All items solved before issue; identities (a)-(b) and the single-photon
%  output of (c) checked numerically by direct quadrature.
% -----------------------------------------------------------------------------
\begin{entry}{Q023}{Fifty per cent loss and the Husimi function}

\entrymeta{
  topic   = {Wigner function; optical loss and thermal noise; Husimi function; nonclassicality},
  est     = {$\sim$40 minutes },
  origin  = Original,
  source  = {Original problem. The converse half of the threshold is proved in
    B.~Nugmanov, N.~Zunikov and F.~Ya.~Khalili, \emph{Robustness of negativity of the
    Wigner function to dissipation}, arXiv:2201.03257 (2022)},
  url     = https://arxiv.org/abs/2201.03257,
  prereq  = {The single-mode Wigner function of \qref{Q022}; Gaussian convolutions},
  status  = drafted,
}

\begin{problem}
\emph{Context.} The ideal hyperedge of \qref{Q023} has a Wigner function that is
negative on part of phase space, and that negativity is what makes such states a
computational resource: if every state in a computation built from Gaussian operations
and homodyne measurements has a non-negative Wigner function, the computation can be
simulated efficiently on a classical computer (A.~Mari and J.~Eisert, Phys.~Rev.~Lett.\
\textbf{109}, 230503 (2012)). Every optical or microwave experiment, however, loses
photons to its environment and picks up thermal noise. This exercise asks how much loss
or noise a negative Wigner function can survive, whatever the input state. The answer
is a universal threshold at $50\%$ transmission. It is one half of the ``$50\%$ rule'';
the other half, that every pure non-Gaussian state stays negative above the threshold,
is proved in the source.

\emph{Setting.} Take one bosonic mode with $[\hat q,\hat p]=i$, $\hbar=1$, write
$z=(q,p)$, and use the Wigner function $W$ of \qref{Q022}, normalised to
$\int dq\,dp\,W=1$. The vacuum has $W_{0}(z)=\pi^{-1}e^{-(q^2+p^2)}$, so each quadrature of
the vacuum has variance $\tfrac12$, the \emph{vacuum variance}. For $v>0$ let
\begin{equation}
  N(x;v):=\frac{1}{\sqrt{2\pi v}}\,e^{-x^{2}/(2v)},\qquad x\in\RR,
  \label{q023:gaussian}
\end{equation}
be the centred Gaussian probability density of variance $v$, and set
$N_v(z):=N(q;v)\,N(p;v)$, so that $W_0=N_{1/2}$.
The convolution of two functions on $\RR^2$ is $(F*G)(z):=\int d^2u\,F(u)\,G(z-u)$, and you
may use $N_a*N_b=N_{a+b}$ (variances of independent Gaussian variables add).

We consider two noise models.
\begin{itemize}
  \item \emph{Pure loss.} The mode enters one port of a beam splitter of power
        transmissivity $T\in(0,1]$, the other port is in the vacuum, and the second
        output is discarded. With $\hat b$ the annihilation operator of the second port,
        the beam splitter is a unitary $U$ with $U\ket{0,0}=\ket{0,0}$ and
        \begin{equation}
          U\hat a^{\dagger}U^{\dagger}=\sqrt T\,\hat a^{\dagger}+\sqrt{1-T}\,\hat b^{\dagger}.
          \label{q023:beam-splitter}
        \end{equation}
  \item \emph{Additive (thermal) noise.} The mode is displaced at random,
        $(q,p)\mapsto(q+\xi_q,\,p+\xi_p)$, with $\xi_q,\xi_p$ independent centred Gaussian
        variables of variance $\bar n$ each. This adds $\bar n$ photons on average.
\end{itemize}
Because the Wigner function transforms under a beam splitter by the corresponding linear
change of phase-space variables, and a partial trace becomes a marginal, both channels act
on $W$ exactly as rescaling plus independent Gaussian noise acts on a probability density:
\begin{equation}
  W_{ out}(q,p)=\int du\,dw\;W(u,w)\,N\bigl(q-\sqrt T\,u;\,v\bigr)\,N\bigl(p-\sqrt T\,w;\,v\bigr),
  \label{q023:channel}
\end{equation}
\begin{equation}
  \text{pure loss: }v=\tfrac{1-T}{2};
  \qquad
  \text{additive noise: }T=1,\ v=\bar n .
  \label{q023:variances}
\end{equation}

The \emph{Husimi function} of a state $\rho$ is
$Q(q,p):=\bra{\alpha}\rho\ket{\alpha}/(2\pi)$, where $\ket\alpha$ is the coherent state
with $\alpha=(q+ip)/\sqrt2$. It is normalised to $\int dq\,dp\,Q=1$ and satisfies $Q\ge0$
for every state, since $\rho\ge0$. A standard fact (K.~E.~Cahill and R.~J.~Glauber,
Phys.~Rev.\ \textbf{177}, 1882 (1969)) relates it to the Wigner function:
\begin{equation}
  Q=W*N_{1/2}.
  \label{q023:husimi}
\end{equation}
A state is \emph{classical} if $\rho=\int d^2\alpha\,P(\alpha)\,\ket\alpha\!\bra\alpha$ with
$P$ a probability density, that is, a statistical mixture of coherent states. Every
classical state has a photon-number variance at least as large as its mean
(L.~Mandel, Opt.~Lett.\ \textbf{4}, 205 (1979)):
\begin{equation}
  \rho\ \text{classical}\ \Longrightarrow\ \operatorname{Var}(\hat n)\ \ge\ \langle\hat n\rangle,
  \qquad \hat n=\hat a^{\dagger}\hat a .
  \label{q023:mandel}
\end{equation}

\begin{enumerate}[label=\textbf{(\alph*)},leftmargin=*]
  \item \emph{Rescaling.} Change variables in \eqref{q023:channel} to show that
        \begin{equation}
          W_{out}(z)=\frac1T\,\bigl[W*N_{v/T}\bigr]\!\Bigl(\frac{z}{\sqrt T}\Bigr),
          \label{q023:rescaled}
        \end{equation}
        that is, the output is the input blurred with variance $v/T$ per quadrature,
        then evaluated at rescaled coordinates.

  \item \emph{The threshold.} Using \eqref{q023:rescaled} and \eqref{q023:husimi}, show
        that additive noise with $\bar n=\tfrac12$ gives $W_{out}=Q$, and that pure loss
        with $T=\tfrac12$ gives $W_{ out}(z)=2\,Q(\sqrt2\,z)$. Then show that
        $W_{out}\ge0$ for every input state whenever $T\le\tfrac12$ or
        $\bar n\ge\tfrac12$. The variance $v/T$ is the added noise referred back to the
        input: how large is it at each threshold, in units of the vacuum variance?

  \item \emph{Positive is not classical.} Let the input be the one-photon state
        $\ket1$. Using \eqref{q023:beam-splitter}, show that the output of pure loss is
        $\rho_T=T\ket1\!\bra1+(1-T)\ket0\!\bra0$, so that by (b) its Wigner function is
        non-negative for $T\le\tfrac12$. Using \eqref{q023:mandel}, show that $\rho_T$ is
        not classical for any $T>0$. What does a non-negative Wigner function therefore
        fail to certify?
\end{enumerate}
\end{problem}

\begin{solution}[Part (a)]
From \eqref{q023:gaussian} we see that $N(x\sqrt{\gamma},v\gamma)=N(x,v)$ for any rescaling factor $\gamma$. So the Wigner function \eqref{q023:channel} can be written as (take $\gamma=1/T$)
\begin{equation}
\begin{split}
     W_{ out}(q,p)&=\int du\,dw\;W(u,w)\,N\bigl(\frac{q}{\sqrt{T}}-\,u;\,\frac{v}{T}\bigr)\,N\bigl(\frac{p}{\sqrt{T}}-\,w;\,\frac{v}{T}\bigr)\\
     &=\int du\,dw\;W(u,w)\,N_{v/T}\left( \frac{q}{\sqrt{T}}-\,u, \frac{p}{\sqrt{T}}-\,w\right)
\end{split}
\end{equation}
which can be written in terms of elements of $\RR^2$, the desired expression clearer by taking $z=\left(\frac{q}{\sqrt{T}}-\frac{p}{\sqrt{T}} \right)$
\begin{equation}
\begin{split}
     W_{ out}(z)&= \int d^2u\, W(u)\,N_{v/T}\left( z-\,u\right)\\
     &= \frac1T\,\bigl[W*N_{v/T}\bigr]\!\Bigl(\frac{z}{\sqrt T}\Bigr)
\end{split}
\end{equation}

\end{solution}

\begin{solution}[Part (b)]
For this additive noise we precisely get that  $W_{out}=W * N_{1/2}$, which is precisely the $Q$ function by \eqref{q023:husimi}. Now, for pure loss with $T=1/2 \implies v= 1/4$, then we get $W_{out}= 2 W * N_{1/2}(\sqrt{2}z)=2Q(\sqrt{2})$. For the additive noise, $\bar n = 1/2 + \delta n$, with $\delta n\geq0$, such that 
\begin{equation}
    W_{out}=W * N_{\bar n} = (W * N_{1/2}) * N_{\delta n},
\end{equation}
and $W * N_{1/2}$ is always non-negative, since it equals the Husimi function at the same time $N_{\delta n}$ is also a non-positive function, so the convolution of both is also non-negative (integral of the product of two non-negative functions)
The analogous arguments follow for $T\leq 1/2$ for the pure loss. $T\leq 1/2 \implies v/T = \frac{1}{2} + \frac{1-2T}{2T} \geq 1/2$ and we can be the same argument for $\frac{1-2T}{2T}\geq0$.
\end{solution}

\begin{solution}[Part (c)]

\end{solution}

% \begin{remark}
% \end{remark}

\begin{takeaway}
\end{takeaway}

%  Everything below is optional.  An entry with no attempt block is complete.
%  Record only what would change what you do next time: typos and arithmetic
%  slips you caught yourself are artefacts of the medium, not facts about the
%  mathematics.  Add "public," to show this block in the public build.
% \attempt{
%   date       = ,               % ISO, YYYY-MM-DD -- required if the block is used
%   time       = ,
%   difficulty = ,               % 1-5, perceived
%   howitwent  = ,
%   insight    = ,
%   note       = ,
% }

\end{entry}